\documentclass[11pt]{article}
\usepackage[margin=1in]{geometry}
\usepackage{amsmath, amssymb}
\usepackage{enumitem}
\usepackage{titlesec}
\usepackage{hyperref}
\usepackage{booktabs}
\usepackage{xcolor}
\usepackage{parskip}

\hypersetup{colorlinks=true, linkcolor=blue, urlcolor=blue}

\title{\textbf{EE627 --- Class Summary}\\[0.4em]
       \large March 6, 2026}
\author{Prof.\ Rensheng Wang}
\date{}

\begin{document}
\maketitle
\tableofcontents
\newpage

%------------------------------------------------------------
\section{Logistic Regression vs.\ Linear Regression}
%------------------------------------------------------------

\subsection{The Fundamental Distinction}
\begin{center}
\begin{tabular}{@{}lll@{}}
\toprule
\textbf{Model} & \textbf{Task} & \textbf{Output} \\
\midrule
Linear regression  & Regression      & Continuous value $\hat{y} \in (-\infty, +\infty)$ \\
Logistic regression & Classification & Probability $p \in [0, 1]$ \\
\bottomrule
\end{tabular}
\end{center}

\subsection{Linear Regression Equation}
\[
\hat{y} = b_0 + b_1 x_1 + b_2 x_2 + \cdots + b_n x_n
\]
The output can be any real number --- unbounded.

\subsection{Logistic Regression Equation}
The same linear combination is fed into the \textbf{sigmoid (logistic) function}
to produce a probability:
\[
\Prob(y = 1) = \frac{1}{1 + \exp\!\bigl(-[b_0 + b_1 x_1 + b_2 x_2 + \cdots + b_n x_n]\bigr)}
\]

\begin{itemize}
    \item Output is always between 0 and 1 --- interpretable as a probability.
    \item Target variable is \textbf{categorical} (class label 0 or 1).
    \item The logistic function squashes the linear combination into
      $[0,1]$ via the S-shaped sigmoid curve.
\end{itemize}

\subsection{Pipeline View}
\[
\underbrace{x_1, x_2, x_3, x_4}_{\text{inputs}}
\;\xrightarrow{\text{linear model}}\;
\underbrace{b_0 + \sum b_i x_i}_{\text{linear score}}
\;\xrightarrow{\text{sigmoid}}\;
\underbrace{p \in [0,1]}_{\text{probability}}
\;\xrightarrow{\text{threshold}}\;
\underbrace{0 \text{ or } 1}_{\text{class label}}
\]


\newpage
%------------------------------------------------------------
\section{Interpreting Logistic Regression Weights: Log Odds}
%------------------------------------------------------------

\subsection{Why Linear Interpretation Breaks Down}
In linear regression, each weight $b_i$ has a direct interpretation: a
one-unit increase in $x_i$ increases $\hat{y}$ by $b_i$. In logistic
regression this is no longer true because the output is a probability,
not a linear function of the inputs.

\subsection{The Log Odds (Logit) Formulation}
To recover a linear interpretation, rearrange the logistic formula:

\[
\log\!\left(\frac{P(y=1)}{1 - P(y=1)}\right)
= \log\!\left(\frac{P(y=1)}{P(y=0)}\right)
= b_0 + b_1 x_1 + b_2 x_2 + \cdots + b_n x_n
\]

\begin{itemize}
    \item The fraction $P(y=1)/P(y=0)$ is called the \textbf{odds} ---
      the probability of the event divided by the probability of no event.
    \item Taking the natural log of the odds gives the \textbf{log odds}
      (also called the \textbf{logit}).
    \item The right-hand side is \emph{linear} in the features.
\end{itemize}

\textit{Key insight:} Logistic regression is a \textbf{linear model for the
log odds}. The weights $b_i$ describe how a one-unit increase in $x_i$
changes the \emph{log odds} of the outcome, not the probability directly.


\newpage
%------------------------------------------------------------
\section{ROC Curve Construction}
%------------------------------------------------------------

The \textbf{Receiver Operating Characteristic (ROC)} curve is the standard
tool for evaluating binary classifiers.

\subsection{Setup}
After fitting a logistic regression model, each sample receives a predicted
probability $p \in [0,1]$. Place the ground-truth labels and predicted
probabilities side by side:

\begin{center}
\begin{tabular}{@{}cc@{}}
\toprule
\textbf{Truth} & \textbf{Predicted prob.} \\
\midrule
1 & 0.9 \\
1 & 0.7 \\
1 & 0.6 \\
1 & 0.4 \\
0 & 0.7 \\
0 & 0.5 \\
0 & 0.4 \\
0 & 0.2 \\
0 & 0.2 \\
\bottomrule
\end{tabular}
\end{center}

\subsection{Threshold Sweep}
Choose a threshold $\tau$ (starting from 1.0, sweeping down to 0):
\[
\hat{y} = \begin{cases} 1 & p \geq \tau \\ 0 & p < \tau \end{cases}
\]

At each threshold, compute:
\begin{itemize}
    \item \textbf{True Positive Rate (TPR):}
      $\displaystyle\frac{\text{correctly predicted positives}}{\text{all actual positives}}$
    \item \textbf{False Positive Rate (FPR):}
      $\displaystyle\frac{\text{negatives incorrectly predicted as positive}}{\text{all actual negatives}}$
\end{itemize}

\subsection{Worked Example}
Using the table above (4 true positives, 5 true negatives):

\begin{center}
\begin{tabular}{@{}lrrrrrrr@{}}
\toprule
\textbf{Threshold} & \textbf{1} & \textbf{0.9} & \textbf{0.7}
  & \textbf{0.6} & \textbf{0.5} & \textbf{0.4} & \textbf{0.2} \\
\midrule
False Positive & 0     & 0     & 1/5   & 1/5   & 2/5   & 3/5   & 1 \\
True Positive  & 0     & 1/4   & 2/4   & 3/4   & 3/4   & 1     & 1 \\
\bottomrule
\end{tabular}
\end{center}

\textit{Reading the table:} at $\tau = 0.7$, the model predicts ``1'' for
any sample with $p \geq 0.7$. The predictions of 1 include truth-1 samples
with prob 0.9 and 0.7 (2 out of 4 true positives captured) and the
truth-0 sample with prob 0.7 (1 false positive out of 5 negatives).

\subsection{Plotting the ROC Curve}
Plot FPR on the $x$-axis and TPR on the $y$-axis. Each threshold setting
gives one point; connecting the points traces the ROC curve.
\begin{itemize}
    \item A \textbf{perfect classifier} reaches the upper-left corner
      $(0, 1)$: zero false positives, all true positives captured.
    \item A \textbf{random classifier} lies on the diagonal from $(0,0)$
      to $(1,1)$.
    \item A better model pushes the curve toward the upper-left.
\end{itemize}


\newpage
%------------------------------------------------------------
\section{AUC: Area Under the ROC Curve}
%------------------------------------------------------------

\subsection{Why AUC?}
Different models produce different ROC curves. To compare them with a
single number, use the \textbf{Area Under the Curve (AUC)}.

\begin{itemize}
    \item AUC ranges from 0 to 1.
    \item Random classifier: $\text{AUC} = 0.5$ (diagonal line).
    \item Perfect classifier: $\text{AUC} = 1.0$.
    \item To improve AUC: push the ROC curve toward the upper-left corner
      (lower FPR at equal TPR, or higher TPR at equal FPR).
\end{itemize}

\subsection{AUC Quality Scale}

\begin{center}
\begin{tabular}{@{}ll@{}}
\toprule
\textbf{AUC range} & \textbf{Test quality} \\
\midrule
0.9 -- 1.0 & Excellent \\
0.8 -- 0.9 & Very good \\
0.7 -- 0.8 & Good \\
0.6 -- 0.7 & Satisfactory \\
0.5 -- 0.6 & Unsatisfactory \\
\bottomrule
\end{tabular}
\end{center}

\subsection{Comparing Multiple Models}
Multiple ROC curves can be plotted on the same axes and compared by their
AUC values. Example: a model with $\text{AUC} = 0.9$ dominates one with
$\text{AUC} = 0.65$ across virtually all threshold settings.


\newpage
%------------------------------------------------------------
\section{R Implementation: Logistic Regression with \texttt{glm}}
%------------------------------------------------------------

In R, logistic regression is fit using \texttt{glm()} with
\texttt{family = "binomial"}.

\subsection{Graduate Admissions Example}
\begin{verbatim}
mydata <- read.csv("binary.csv")
mylogit <- glm(admit ~ gre + gpa + rank,
               data = mydata, family = "binomial")
summary(mylogit)
\end{verbatim}

Abbreviated output:
\begin{verbatim}
Coefficients:
            Estimate Std. Error z value Pr(>|z|)
(Intercept) -3.98998    1.13995   -3.50  0.00047 ***
gre          0.00226    0.00109    2.07  0.03847 *
gpa          0.80404    0.33182    2.42  0.01539 *
rank2       -0.67544    0.31649   -2.13  0.03283 *
rank3       -1.34020    0.34531   -3.88  0.00010 ***
rank4       -1.55146    0.41783   -3.71  0.00020 ***

Null deviance: 499.98  on 399 degrees of freedom
Residual deviance: 458.52  on 394 degrees of freedom
AIC: 470.5
\end{verbatim}

\subsection{Interpreting the Coefficients}
Because logistic regression models log odds, each coefficient
describes the change in log odds per unit increase in the predictor:
\begin{itemize}
    \item \textbf{GRE:} one-unit increase $\Rightarrow$ log odds of
      admission increase by $+0.002$.
    \item \textbf{GPA:} one-unit increase $\Rightarrow$ log odds
      increase by $+0.804$.
    \item \textbf{Rank 2 vs.\ Rank 1:} attending a rank-2 institution
      (vs.\ rank-1) changes log odds by $-0.675$ (lower chance of
      admission).
    \item Rank 3 and Rank 4 have progressively more negative
      coefficients, confirming a monotone penalty.
\end{itemize}

\subsection{ROC Curve in R}
The ROC curve can be constructed manually or via the \texttt{ROCR}
package:
\begin{verbatim}
library(ROCR)
p   <- predict(model, newdata=test, type="response")
pr  <- prediction(p, test$Survived)
prf <- performance(pr, measure="tpr", x.measure="fpr")
plot(prf)
auc <- performance(pr, measure="auc")@y.values[[1]]
\end{verbatim}


\newpage
%------------------------------------------------------------
\section{Case Study: Titanic Survival Prediction}
%------------------------------------------------------------

The Titanic dataset is a standard binary classification benchmark.
The task is to predict survival (0/1) from passenger features.

\subsection{Features Used}
From the raw CSV, columns selected:
\begin{itemize}
    \item \textbf{Survived} (target), \textbf{Pclass}, \textbf{Sex},
      \textbf{Age}, \textbf{SibSp}, \textbf{Parch}, \textbf{Fare},
      \textbf{Embarked}
\end{itemize}

\subsection{Missing Value Handling}
\begin{itemize}
    \item \textbf{Age:} missing values replaced with the column mean.
    \item \textbf{Embarked:} rows with missing values dropped entirely
      (very few affected).
\end{itemize}

\subsection{Workflow}
\begin{verbatim}
data <- subset(raw, select=c(2,3,5,6,7,8,10,12))
data$Age[is.na(data$Age)] <- mean(data$Age, na.rm=TRUE)
data <- data[!is.na(data$Embarked),]

train <- data[1:800,]
test  <- data[801:889,]

model <- glm(Survived ~., family=binomial(link='logit'),
             data=train)

fitted <- predict(model, newdata=test, type='response')
fitted <- ifelse(fitted > 0.5, 1, 0)
accuracy <- 1 - mean(fitted != test$Survived)
\end{verbatim}

\subsection{Key Steps}
\begin{enumerate}
    \item Load data and inspect missingness (\texttt{missmap}).
    \item Impute / drop missing values.
    \item Fit \texttt{glm} with \texttt{binomial} family on training set.
    \item Predict probabilities on test set; apply threshold 0.5.
    \item Evaluate: print accuracy; plot ROC; compute AUC.
\end{enumerate}

\end{document}
